% GENERATED FILE. Edit canonical lessons, then run npm run generate:books.
% canonical-source-hash: fnv1a64:8c64d672cc1175cc
% canonical-lessons: SA-C264-karisyami, SA-C264-pathisyami, SA-C264-near-term-plan, SA-C264-mixed-plan-practice, SA-C264-first-person-future-retrieval

\chapter{What I will do next}
\label{ch:sa-near-term-first-person-plans}
\hlchaptermodality{264}

\begin{chapteropening}
\textbf{By the end of this chapter:} \emph{I can make a near-term first-person plan with three taught future forms.}

Complete chapter 264 by independently choosing a time and a taught first-person future action.
\end{chapteropening}

\section[\texorpdfstring{\sk{करिष्यामि} (kariṣyāmi)}{करिष्यामि (kariṣyāmi)}]{\sk{करिष्यामि} (kariṣyāmi) — I will do}

\label{lesson:SA-C264-karisyami}



\begin{quote}
Say “I do,” \textbf{\sk{करोमि}}, then “I will go,” \textbf{\sk{गमिष्यामि}}.
\end{quote}

\subsection*{You'll want to know: \sk{करिष्यामि}}
\textbf{\sk{करिष्यामि}} — \emph{kariṣyāmi} — means “I will do” or “I will make.” Keep the whole future form together: \textbf{\sk{करिष्यामि}}.

\begin{cousinweb}
\textbf{\sk{करिष्यामि}} comes from the root \textbf{\sk{कृ}} (\emph{kṛ}), “do, make,” already met inside \textbf{\sk{करोमि}}. The future stem appears as \textbf{\sk{करिष्य}-} (\emph{kariṣya-}), followed here by first-person \textbf{-\sk{आमि}}. Whitney, \emph{Sanskrit Grammar} §§932–935, describes the Sanskrit s-future and its \textbf{-iṣya- / -sya-} formation.
\end{cousinweb}

\subsection*{Your turn}
Say these aloud:

\begin{itemize}
\raggedright
  \item “I will do” — \emph{kariṣyāmi}
  \item “Tomorrow, I will do the work” — \emph{śvaḥ kāryaṃ kariṣyāmi}
  \item choose “I will go” — \emph{gamiṣyāmi}; now “I will do” — \emph{kariṣyāmi}
\end{itemize}

\begin{tcolorbox}[breakable,colback=teal!4,colframe=teal!35!black,title={Before you move on}]
Say “I will do.” (\textbf{\sk{करिष्यामि}}.)
\end{tcolorbox}
\section[\texorpdfstring{\sk{पठिष्यामि} (paṭhiṣyāmi)}{पठिष्यामि (paṭhiṣyāmi)}]{\sk{पठिष्यामि} (paṭhiṣyāmi) — I will read}

\label{lesson:SA-C264-pathisyami}



\begin{quote}
Retrieve “I will go” and “I will do.” (\emph{Gamiṣyāmi; kariṣyāmi.})
\end{quote}

\subsection*{You'll want to know: \sk{पठिष्यामि}}
\textbf{\sk{पठिष्यामि}} — \emph{paṭhiṣyāmi} — means “I will read” or “I will study.” Learn the complete form; its future stem is given, not guessed.

\begin{cousinweb}
\textbf{\sk{पठिष्यामि}} comes from the root \textbf{\sk{पठ्}} (\emph{paṭh}), “read, recite.” Its future here is \textbf{\sk{पठिष्य}-} (\emph{paṭhiṣya-}), followed by first-person \textbf{-\sk{आमि}}. Whitney, \emph{Sanskrit Grammar} §§932–935, describes the Sanskrit s-future and the \textbf{-iṣya- / -sya-} formation used here.
\end{cousinweb}

\subsection*{Your turn}
Say these aloud:

\begin{itemize}
\raggedright
  \item “I will read” — \emph{paṭhiṣyāmi}
  \item “Tomorrow, I will read” — \emph{śvaḥ paṭhiṣyāmi}
  \item “I will do; I will read” — \emph{kariṣyāmi; paṭhiṣyāmi}
\end{itemize}

\begin{tcolorbox}[breakable,colback=teal!4,colframe=teal!35!black,title={Before you move on}]
Say “I will read.” (\textbf{\sk{पठिष्यामि}}.)
\end{tcolorbox}
\section[(near-term first-person plan)]{A small plan: when, then what I will do}

\label{lesson:SA-C264-near-term-plan}



\begin{quote}
Recall tomorrow, the day after tomorrow, and the three complete forms “I will go,” “I will do,” and “I will read.”
\end{quote}

\begin{grammarlens}[title={Grammar Lens: time first, plan second}]
Make a near-term plan in two small choices: choose \textbf{\sk{श्वः}} (\emph{śvaḥ}) or \textbf{\sk{परश्वः}} (\emph{paraśvaḥ}), then choose one complete first-person future form.
\end{grammarlens}

\subsection*{Your turn}
Say these aloud:

\begin{itemize}
\raggedright
  \item “Tomorrow, I will read” — \emph{śvaḥ paṭhiṣyāmi}
  \item “The day after tomorrow, I will go” — \emph{paraśvaḥ gamiṣyāmi}
  \item “Tomorrow, I will do the work” — \emph{śvaḥ kāryaṃ kariṣyāmi}
\end{itemize}

\begin{tcolorbox}[breakable,colback=teal!4,colframe=teal!35!black,title={Before you move on}]
What are the two choices in a small plan? (When; then what I will do.)
\end{tcolorbox}
\section[Practice]{Three plans, no copied model}

\label{lesson:SA-C264-mixed-plan-practice}



\begin{quote}
Keep the speaker ending in every answer: this round is all “I.”
\end{quote}

\subsection*{Your turn}
There is no Sanskrit model to copy. Build each answer from the English cue.

Say these aloud:

\begin{itemize}
\raggedright
  \item tomorrow + I will go — \emph{śvaḥ gamiṣyāmi}
  \item tomorrow + I will read — \emph{śvaḥ paṭhiṣyāmi}
  \item day after tomorrow + I will do the work — \emph{paraśvaḥ kāryaṃ kariṣyāmi}
\end{itemize}

\begin{tcolorbox}[breakable,colback=teal!4,colframe=teal!35!black,title={Before you move on}]
Give the three future forms only: I will go, I will do, I will read. (\emph{Gamiṣyāmi; kariṣyāmi; paṭhiṣyāmi.})
\end{tcolorbox}
\section[Practice]{Say what you will do}

\label{lesson:SA-C264-first-person-future-retrieval}



\begin{quote}
This is a \textbf{first-person singular} checkpoint. Every answer tells your own near-term plan.
\end{quote}

\subsection*{Your turn — independent retrieval}
Say these aloud:

\begin{itemize}
\raggedright
  \item You plan to travel tomorrow. Say when and what you will do. — \emph{śvaḥ gamiṣyāmi}
  \item You plan to study tomorrow. Say when and what you will do. — \emph{śvaḥ paṭhiṣyāmi}
  \item You plan to do the work the day after tomorrow. — \emph{paraśvaḥ kāryaṃ kariṣyāmi}
  \item Choose a different one of the three actions and make your own tomorrow plan.
\end{itemize}

\begin{grammarlens}[title={Grammar Lens: the boundary}]
These three complete forms make a useful first-person singular ramp. This is \textbf{not a complete future paradigm}. Second person, dual, and plural forms, and the future stems of other roots, wait for later lessons.
\end{grammarlens}

\begin{tcolorbox}[breakable,colback=teal!4,colframe=teal!35!black,title={Before you move on}]
Without a model, say one true or imagined plan for tomorrow with \textbf{\sk{गमिष्यामि}}, \textbf{\sk{करिष्यामि}}, or \textbf{\sk{पठिष्यामि}}.
\end{tcolorbox}
